% Section A.4, the solutions to the self-assessment exercises in Chapter 10: Sets 4 to 6, from
% lectures/L10/appendix/c_exercises.md, and Exercise 10.18, from the exercise of
% lectures/L10/appendix/d_embedded_constraints.md. The repository has no solutions to these; they
% are written for the book.

\section{Chapter 10: The Flatten Layer and the Complete Network}
\label{sol:sec:c10}

\solution{c10:ex:10}{Six feedforward algorithms}
What gives each layer away is its parameters and the shape of its loop: what is summed, over what,
and whether the same weights come back at every position.

\step{A: a linear regression layer}
One input, one weight $k$, one bias $m$, and no activation function: $y = kx + m$, exactly
\code{predict()} in Chapter~\ref{c1:ch:linreg}.

\step{B: a dense layer}
Every node \code{n} has its own bias and its own weight for every input, the weighted sum over
\emph{all} inputs starts at the bias, and the activation function is applied to it: the
feedforward of \secref{c5:sec:math}.

\step{C: a conv layer}
The loop runs over output positions, and each output is the weighted sum of a small
\code{kernelSize × kernelSize} window of the input, starting at \code{(row, col)}, plus a bias,
passed through the activation function. Two things separate it from B: the sum only covers a local
window (locality), and \code{kernel} and \code{bias} carry no index for the output position, so
the same weights are used at every position (weight sharing). The \code{input} it reads is the
zero-padded one, which is why \code{input[row + kr][col + kc]} never runs off its edge.

\step{D: a max pooling layer}
Each output is the largest value in a non-overlapping \code{poolSize × poolSize} block: the block
steps by \code{poolSize}, so the stride equals the pool size. There is no weight or bias anywhere,
and nothing is summed.

\step{E: a flatten layer}
Nothing is computed at all; the data is only given a new shape, from \code{height × width} to one
row of \code{height * width} values.

\step{F: a dense layer with one node and no activation function}
It has the dense layer's structure, a bias and one weight per input summed over all the inputs,
but only for node 0, and with the activation function left out (\code{ActFunc::None}). That makes
it linear regression generalised to several inputs, $y = m + \sum_i k_i x_i$: with a single input
it is A written out as a loop. It isn't A as written, because it takes a vector of inputs rather
than one, and it isn't B in general, because B has any number of nodes and applies an activation
function. \secref{c10:sec:course} makes the same point from the other side: \code{Fixed} is a
\code{Dense} layer with one node and no activation function. Either name earns the answer, as long
as the reason is given.

\solution{c10:ex:11}{Tracing shapes}
\begin{narrowtable}{@{}ll@{}}
  \thead{After} & \thead{Shape} \\
  \midrule
  conv layer & 8\,×\,8 \\
  max pooling layer & 4\,×\,4 \\
  flatten layer & 1\,×\,16 \\
  dense layer & 1\,×\,10 \\
\end{narrowtable}

\begin{itemize}
  \item \textbf{Conv.} The padding is \code{kernelSize / 2} $= 3 / 2 = 1$ on each side, so with
    stride 1 the output is $\lfloor (8 + 2 \cdot 1 - 3) / 1 \rfloor + 1 = 8$ wide: the same size as
    the input, which is what the padding is for.
  \item \textbf{Max pooling.} Non-overlapping 2\,×\,2 blocks halve each side: $8 / 2 = 4$.
  \item \textbf{Flatten.} $4 \cdot 4 = 16$ values in one row.
  \item \textbf{Dense.} One output per node: 10 values.
\end{itemize}

\solution{c10:ex:12}{Counting trainable parameters}
\begin{narrowtable}{@{}lr@{}}
  \thead{Layer} & \thead{Trainable parameters} \\
  \midrule
  conv & $3 \cdot 3 + 1 = 10$ \\
  max pooling & 0 \\
  flatten & 0 \\
  dense & $10 \cdot 16 + 10 = 170$ \\
  \midrule
  total & 180 \\
\end{narrowtable}

The conv layer has one 3\,×\,3 kernel and one bias, whatever the size of the image. The dense layer
has 10 nodes, each with one weight per input (the 16 values from the flatten layer) and a bias of
its own. Max pooling and flatten are purely structural: one picks values that are already there,
the other changes their shape, and neither has anything to learn. More than 94\% of the network's
parameters are in the dense layer at the end.

\solution{c10:ex:13}{What max pooling has to remember}
It has to remember \textbf{which position in each block held the maximum}, or enough to find it
again. The standalone \code{ml::MaxPoolLayer} of Chapter~\ref{c7:ch:layers} stores the index of
each block's maximum; \code{MaxPool} in Chapter~\ref{c9:ch:maxpool} stores the input itself,
\code{myInput}, and searches each block again.

The reason is how the gradient is routed. Only the position that produced a pooled value
influenced the output, so \code{backpropagate()} must hand that position the whole incoming
gradient and every other position in the block zero. The output alone can't say where the maximum
was: a block holding \code{1.9} could have held it in any of its four cells. And the search in
\code{backpropagate()} has to break ties exactly as \code{feedforward()} did, towards the first
occurrence in row-major order, or the gradient lands on a pixel that didn't produce the output.

\solution{c10:ex:14}{Conv vs.\ dense, parameter count}
\textbf{10}: nine kernel weights and one bias, against the dense layer's 65.

The difference comes from the two constraints a conv layer adds to a dense one:
\begin{itemize}
  \item \textbf{Locality.} Each output looks at a 3\,×\,3 window, 9 pixels, rather than at all 64.
  \item \textbf{Weight sharing.} The same kernel is reused at every one of the 64 positions,
    instead of every output having weights of its own.
\end{itemize}
The comparison even flatters the dense layer. Its 65 parameters produce \emph{one} output; the conv
layer's 10 produce 64 of them, one per position. A dense layer producing the same 64 outputs from
the same 64 inputs would need $64 \cdot 64 + 64 = 4160$ parameters. And the gap grows with the
image: the dense layer's count grows with the number of pixels, while the conv layer's stays at 10
for any image size, because it only depends on the kernel.

\solution{c10:ex:15}{Flatten in reverse}
Back into a \textbf{3\,×\,4} matrix: the shape of the input the layer received, not of the output
it produced.

\code{backpropagate()} computes the gradients for the layer's \emph{input}, and hands them to the
layer before it, which produced a 3\,×\,4 matrix and expects exactly one gradient per value it
output. The reshape also has to be the exact inverse of the one in \code{feedforward()}: in
row-major order, element $(i, j)$ went to index $4i + j$ of the vector, so the gradient at index
$4i + j$ has to come back to $(i, j)$, the position whose value it belongs to. Any other order puts
each gradient on a value it has nothing to do with, and the layers behind learn from noise.

The flatten layer in Chapter~\ref{c10:ch:flatten} only accepts square inputs, but the rule is the
same for any shape.

\solution{c10:ex:16}{True or false?}
\begin{enumerate}
  \item \textbf{False.} A max pooling layer has no trainable parameters at all: it only selects the
    largest value in each block. Its \code{optimize()} exists because the interface it shares with
    \code{Conv} declares one, and it's a genuine no-op that returns \code{true}.
  \item \textbf{True.} That is weight sharing: one kernel slides across the whole image, which is
    what keeps the parameter count down and lets the same pattern be recognized wherever it
    appears.
  \item \textbf{False.} It performs the \emph{inverse} reshape: \code{feedforward()} goes from 2D
    to 1D, and \code{backpropagate()} goes from 1D back to 2D, so each gradient lands on the value
    it came from.
  \item \textbf{True}, for a single input. Both predict $y = wx + b$, and both update with the
    same rule: with no activation function the derivative is 1, so $\Delta e = \delta$, and
    $b$ moves by $\delta \cdot \LR$ and $w$ by $\delta \cdot \LR \cdot x$, exactly as $m$ and $k$
    do in Chapter~\ref{c1:ch:linreg}. Two differences in the book's implementations don't change
    the model, only where training starts and one shortcut: \code{ml::lin_reg::Fixed} starts at
    zero while \code{Dense} draws random start values, and \code{Fixed::optimize()} sets the bias
    directly when $x = 0$ (Exercise~\ref{c1:ex:13}). With several inputs, the dense node is linear
    regression over several inputs.
  \item \textbf{True.} Each extra kernel adds its own weights (and, usually, a bias of its own), so
    the parameter count grows. Each kernel produces its own feature map of the same width and
    height, so the output gains maps, or channels, not size; width and height are set by the input
    size, the kernel size, the padding and the stride. This book only ever uses one kernel per conv
    layer, to keep things simple.
\end{enumerate}

\solution{c10:ex:17}{Short answer}
\begin{enumerate}[start=6]
  \item A dense layer gives every output its own weight for every input pixel, so its parameter
    count grows with the size of the image; a conv layer combines only a small local window
    (locality) and reuses the same kernel at every position (weight sharing), so its parameter
    count is fixed by the kernel size alone.
  \item \textbf{Incorrect.} A max pooling layer has no weights at all, so there is nothing for it
    to update; its \code{optimize()} does nothing. If anything it makes training \emph{cheaper}: it
    shrinks the data every layer after it has to process, a 2\,×\,2 pool by a factor of four, which
    means fewer inputs, and fewer weights, in the dense layer behind it.
  \item \textbf{\code{MaxPool::optimize()}} is a genuine no-op, not a placeholder: it ignores the
    learning rate and returns \code{true}. It has to exist because \code{MaxPool} shares
    \code{conv_layer::Interface} with \code{Conv}, which does need \code{optimize()}, and sharing
    the interface is what lets the network hold both in one list and call the same three methods
    on each.

    \textbf{\code{flatten_layer::Interface}} belongs to flatten layers alone, and a flatten layer
    can never have trainable parameters, so there is nothing an \code{optimize()} could ever do.
    Declaring it and leaving it empty would promise something no flatten layer does, and would
    leave it to whoever writes the class to remember that the method must do nothing. Leaving it
    out moves that fact into the type system: the network simply doesn't call \code{optimize()} on
    its flatten layer, and a \code{Flatten} that tried to optimize something wouldn't compile.
\end{enumerate}

\solution{c10:ex:18}{Memory footprint}
\step{The network from Chapters 4 and 5}
Chapter~\ref{c5:ch:dense} trains a hidden layer of 8 nodes over the 2 inputs of the XOR pattern,
and an output layer of 1 node over those 8:
\begin{align*}
  \text{hidden layer} &= 8 \cdot 2 \text{ weights} + 8 \text{ biases} = 24 \text{ values} \\
  \text{output layer} &= 1 \cdot 8 \text{ weights} + 1 \text{ bias} = 9 \text{ values} \\
  \text{total} &= 33 \text{ values}
\end{align*}

\begin{narrowtable}{@{}lrr@{}}
  \thead{Type} & \thead{Bytes each} & \thead{Footprint} \\
  \midrule
  \code{double} & 8 & $33 \cdot 8 = 264$ B \\
  \code{float} & 4 & $33 \cdot 4 = 132$ B \\
  \code{int8} & 1 & $33 \cdot 1 = 33$ B \\
\end{narrowtable}

\noindent The smaller network of Chapter~\ref{c4:ch:network}, 3 hidden nodes over 2 inputs and 1
output node over 3, is the one \secref{c10:sec:memory} counts: 13 values, so 104, 52 and 13 bytes.

\step{Against the layer sized for 28 × 28 images}
The course's networks fit comfortably at any of the three precisions on any target: even as
\code{double}, a few hundred bytes. The layer of 128 nodes over 784 inputs is another matter, with
64\,KB $= 65{,}536$ bytes and 512\,KB $= 524{,}288$ bytes:
\begin{itemize}
  \item \textbf{64\,KB of RAM: none of the three fits.} Even at \code{int8}, the layer's
    100,480 bytes (98.125\,KB) are more than the whole RAM.
  \item \textbf{512\,KB of RAM: \code{float} and \code{int8} fit, \code{double} does not.}
    401,920 bytes (392.5\,KB) and 100,480 bytes fit; 803,840 bytes (785\,KB) do not.
\end{itemize}
``Plausibly'' is the right word, because the RAM has to hold more than the weights: the layers'
outputs, errors and gradients, the stack, and the rest of the program's data. \code{float} in
512\,KB leaves about 120\,KB for all of that, and a real network has more than one layer. This is
the argument for training at full precision and deploying at \code{int8}, from
\secref{c10:sec:quantization}.
